% titel: Blatt 9, A11: Löslichkeitsprodukt Ag2CrO4 und AgBr
% skript: Mortimer Kap. 15 (Löslichkeitsprodukt)

% Teil I: Werkzeuge, kein Lösungsweg.

Das Löslichkeitsprodukt L verknüpft die Ionenkonzentrationen einer
schwerlöslichen Verbindung mit ihrer molaren Löslichkeit.

\begin{werkzeugkasten}[Löslichkeitsprodukt]
Für $A_aB_b\rightleftharpoons aA^{m+}+bB^{n-}$:
\[
  L = [A^{m+}]^a\,[B^{n-}]^b .
\]
Bei reiner Lösung einer Substanz gilt $[A]=a s$, $[B]=b s$, wobei $s$ die
molare Löslichkeit ist. Für Ag$_2$CrO$_4$:
\[
  \mathrm{Ag_2CrO_4 \rightleftharpoons 2 Ag^+ + CrO_4^{2-}},\qquad
  L = (2s)^2\cdot s = 4s^3.
\]
\textit{Skript:} Mortimer Kap. 15 (Löslichkeitsprodukt).
\end{werkzeugkasten}

\begin{achtung}[Darauf achten]
\begin{itemize}
  \item Bei Fremdionenzusatz (gemeinsames Ion, hier Br$^-$ aus NaBr) gilt
        $L=[\text{M}^{n+}][\text{Anion}_{\text{ges}}]$; die Löslichkeit $s$
        fällt entsprechend.
  \item Exponenten: $L(AgBr)=[Ag^+][Br^-]$ (1:1), $L(Ag_2CrO_4)=4s^3$.
  \item Einheiten: $L$ trägt die Einheit mol$^2$L$^{-2}$ (AgBr) bzw.
        mol$^3$L$^{-3}$ (Ag$_2$CrO$_4$).
  \item Das Ionenprodukt Q gegen L vergleichen (Fällung, falls $Q>L$).
\end{itemize}
\end{achtung}
